




\documentclass{ximera}

\input{../../../preamble.tex}


\definecolor{fvdnavy}{HTML}{071D43}
\definecolor{fvdcyan}{HTML}{20BFC6}
\definecolor{fvdcyanlight}{HTML}{EFFCFB}
\definecolor{fvdmagenta}{HTML}{F10D69}
\definecolor{fvdmagentalight}{HTML}{FFF1F7}
\definecolor{fvdtext}{HTML}{163044}





%=================================================
% Pedagogische omgevingen
% PDF: tcolorbox
% HTML: learning-box
%=================================================

\newenvironment{voorkennis}
{%
    \ifdefined\HCode
        \HCode{%
<section class="learning-box learning-box--prior">
<div class="learning-box__header">
<span class="learning-box__icon" aria-hidden="true"></span>
<span class="learning-box__title">Voorkennis</span>
</div>
<div class="learning-box__content">
        }%
        \begin{itemize}
    \else
       \begin{tcolorbox}[
    enhanced,
    breakable,
    colback=fvdcyanlight,
    colframe=fvdcyan,
    coltext=fvdtext,
    boxrule=0.5pt,
    leftrule=4pt,
    arc=4mm,
    outer arc=4mm,
    left=7mm,
    right=6mm,
    top=5mm,
    bottom=5mm,
    before skip=6mm,
    after skip=6mm,
    title=Voorkennis,
    coltitle=fvdcyan,
    fonttitle=\bfseries\Large,
    attach title to upper={\par\smallskip},
    boxed title style={
        empty,
        left=0mm,
        right=0mm,
        top=0mm,
        bottom=0mm
    },
    overlay={
        \draw[
            fvdcyan,
            line width=0.5pt
        ]
        ([xshift=42mm,yshift=-13mm]frame.north west)
        --
        ([xshift=-7mm,yshift=-13mm]frame.north east);
    }
]
        \begin{itemize}
    \fi
}
{%
    \end{itemize}
    \ifdefined\HCode
        \HCode{%
</div>
</section>
        }%
    \else
        \end{tcolorbox}
    \fi
}


\newenvironment{leerdoelen}
{%
    \ifdefined\HCode
        \HCode{%
<section class="learning-box learning-box--goals">
<div class="learning-box__header">
<span class="learning-box__icon" aria-hidden="true"></span>
<span class="learning-box__title">Leerdoelen</span>
</div>
<div class="learning-box__content">
        }%
        \begin{itemize}
    \else
\begin{tcolorbox}[
    enhanced,
    breakable,
    colback=fvdmagentalight,
    colframe=fvdmagenta,
    coltext=fvdtext,
    boxrule=0.5pt,
    leftrule=4pt,
    arc=4mm,
    outer arc=4mm,
    left=7mm,
    right=6mm,
    top=5mm,
    bottom=5mm,
    before skip=6mm,
    after skip=6mm,
    title=Leerdoelen,
    coltitle=fvdmagenta,
    fonttitle=\bfseries\Large,
    attach title to upper={\par\smallskip},
    boxed title style={
        empty,
        left=0mm,
        right=0mm,
        top=0mm,
        bottom=0mm
    },
    overlay={
        \draw[
            fvdmagenta,
            line width=0.5pt
        ]
        ([xshift=42mm,yshift=-13mm]frame.north west)
        --
        ([xshift=-7mm,yshift=-13mm]frame.north east);
    }
]
        \begin{itemize}
    \fi
}
{%
    \end{itemize}
    \ifdefined\HCode
        \HCode{%
</div>
</section>
        }%
    \else
        \end{tcolorbox}
    \fi
}


\addPrintStyle{../../..}

\begin{document}

% FVD_CHAPTER_DATA_START
% Automatisch gegenereerd uit index.tex.
% Niet handmatig aanpassen.
\fvdchapterdata{3}{Afleiden van samengestelde functie}{1}
% FVD_CHAPTER_DATA_END










\title{Producten en quotiënten afleiden}

\begin{abstract}
Tot nu toe kon je vooral functies afleiden die geschreven zijn als een som
van eenvoudige termen.

Maar functies komen vaak voor als een product of een quotiënt.
Denk bijvoorbeeld aan
\[
f(x)=(x^2+1)(3x-2)
\qquad\text{of}\qquad
g(x)=\frac{x^2+1}{x-2}.
\]

In dit hoofdstuk leer je hoe je zulke functies afleidt.
Daarbij is niet alleen het correct toepassen van een formule belangrijk:
je leert ook herkennen welke afleidingsregel geschikt is en wanneer
herschrijven eenvoudiger is.

Die keuzes worden later essentieel bij het onderzoeken van rationale
functies en bij het oplossen van limietproblemen.
\end{abstract}

\maketitle


% ============================================================
% VOORKENNIS EN LEERDOELEN
% ============================================================

\begin{voorkennis}
  \item Je kan veeltermfuncties afleiden.
  \item Je kent de machtsregel
        \[
        (x^n)'=nx^{n-1}.
        \]
  \item Je kan de som- en verschilregel gebruiken.
  \item Je kan machten met negatieve exponenten vereenvoudigen.
  \item Je kent het domein van een rationale functie.
  \item Je kan algebraïsche uitdrukkingen vereenvoudigen en ontbinden.
\end{voorkennis}

\begin{leerdoelen}
  \item Je kunt de productregel gebruiken om een product van functies af te leiden.
  \item Je kunt de quotiëntregel gebruiken om een quotiënt van functies af te leiden.
  \item Je kunt uitleggen waarom
        \[
        (uv)'\neq u'v'
        \qquad\text{en}\qquad
        \left(\frac{u}{v}\right)'\neq\frac{u'}{v'}.
        \]
  \item Je kunt in een functie de factoren of de teller en noemer herkennen.
  \item Je kunt bepalen welke afleidingsregel geschikt is.
  \item Je kunt herkennen wanneer herschrijven eenvoudiger is dan het toepassen
        van de product- of quotiëntregel.
  \item Je kunt verschillende afleidingsregels combineren.
  \item Je kunt een gevonden afgeleide algebraïsch vereenvoudigen.
  \item Je kunt bij rationale functies rekening houden met het domein.
\end{leerdoelen}

% ============================================================
% 1. NIET ELKE BEWERKING MAG JE ZOMAAR AFLEIDEN
% ============================================================

\FVDsection{Niet elke bewerking mag je zomaar afleiden}

Voor sommen en verschillen ken je eenvoudige regels:
\[
(u+v)'=u'+v'
\]
en
\[
(u-v)'=u'-v'.
\]

Het is verleidelijk om te denken dat dit voor elke bewerking op dezelfde
manier werkt.

Bij een product zou je bijvoorbeeld kunnen vermoeden dat
\[
(uv)'=u'v'.
\]

Dat is echter niet juist.

\begin{voorbeeld}
Neem
\[
u(x)=x
\qquad\text{en}\qquad
v(x)=x.
\]

Dan is
\[
u(x)v(x)=x^2,
\]
dus
\[
(uv)'(x)=2x.
\]

Maar
\[
u'(x)v'(x)=1\cdot1=1.
\]

Bijgevolg:
\[
(uv)'\neq u'v'.
\]
\end{voorbeeld}

Hetzelfde probleem ontstaat bij een quotiënt. In het algemeen geldt ook
niet
\[
\left(\frac{u}{v}\right)'=\frac{u'}{v'}.
\]

Voor producten en quotiënten hebben we dus nieuwe afleidingsregels nodig.


% ============================================================
% 2. DE PRODUCTREGEL
% ============================================================

\FVDsection{De productregel}

Beschouw twee afleidbare functies $u$ en $v$ en hun product
\[
f(x)=u(x)v(x).
\]

Wanneer $x$ verandert, kunnen \emph{beide} factoren veranderen.
De verandering van het product wordt daarom niet alleen bepaald door
$u'$ of alleen door $v'$.

\begin{eigenschap}[Productregel]
Als $u$ en $v$ afleidbaar zijn, dan geldt
\[
\boxed{
(uv)'=u'v+uv'
}
\]

of, volledig uitgeschreven,
\[
\boxed{
\frac{d}{dx}\bigl(u(x)v(x)\bigr)
=
u'(x)v(x)+u(x)v'(x).
}
\]
\end{eigenschap}

Een handige geheugensteun is:

\[
\boxed{
\text{afgeleide eerste}\cdot\text{tweede}
+
\text{eerste}\cdot\text{afgeleide tweede}
}
\]


% ------------------------------------------------------------
% Voorbeeld 1
% ------------------------------------------------------------

\begin{voorbeeld}
Bepaal de afgeleide van
\[
f(x)=(x^2+1)(3x-2).
\]

We herkennen twee factoren:
\[
u(x)=x^2+1,
\qquad
v(x)=3x-2.
\]

Hun afgeleiden zijn
\[
u'(x)=2x,
\qquad
v'(x)=3.
\]

Met de productregel:
\[
f'(x)
=
u'(x)v(x)+u(x)v'(x).
\]

Dus
\[
f'(x)
=
2x(3x-2)+(x^2+1)\cdot3.
\]

We vereenvoudigen:
\[
\begin{aligned}
f'(x)
&=6x^2-4x+3x^2+3\\
&=9x^2-4x+3.
\end{aligned}
\]

Bijgevolg:
\[
\boxed{f'(x)=9x^2-4x+3.}
\]
\end{voorbeeld}


% ============================================================
% 3. MOET JE ALTIJD DE PRODUCTREGEL GEBRUIKEN?
% ============================================================

\FVDsection{Moet je altijd de productregel gebruiken?}

Een formule herkennen betekent nog niet dat je ze automatisch moet
gebruiken.

Soms is het slimmer om de functie eerst te herschrijven.

\begin{voorbeeld}
Bepaal de afgeleide van
\[
f(x)=x^2(x+3).
\]

\textbf{Methode 1: productregel}

\[
\begin{aligned}
f'(x)
&=2x(x+3)+x^2\\
&=2x^2+6x+x^2\\
&=3x^2+6x.
\end{aligned}
\]

\textbf{Methode 2: eerst uitwerken}

Omdat
\[
f(x)=x^3+3x^2,
\]
krijgen we onmiddellijk
\[
f'(x)=3x^2+6x.
\]

Beide methoden geven uiteraard hetzelfde resultaat.
\end{voorbeeld}

\begin{opmerking}
Een goede strategie is niet noodzakelijk de strategie met de
meeste formules.

Vraag jezelf vóór je begint af:

\[
\boxed{\text{Kan ik de functie eerst eenvoudiger schrijven?}}
\]
\end{opmerking}


% ============================================================
% 4. DE QUOTIËNTREGEL
% ============================================================

\FVDsection{De quotiëntregel}

We bekijken nu een functie van de vorm
\[
f(x)=\frac{u(x)}{v(x)},
\qquad
v(x)\neq0.
\]

Ook hier veranderen teller en noemer wanneer $x$ verandert.

De afgeleide is \emph{niet}
\[
\frac{u'}{v'}.
\]

Daarvoor geldt een aparte regel.

\begin{eigenschap}[Quotiëntregel]
Als $u$ en $v$ afleidbaar zijn en $v(x)\neq0$, dan geldt
\[
\boxed{
\left(\frac{u}{v}\right)'
=
\frac{u'v-uv'}{v^2}
}
\]

of volledig uitgeschreven:
\[
\boxed{
\frac{d}{dx}
\left(
\frac{u(x)}{v(x)}
\right)
=
\frac{
u'(x)v(x)-u(x)v'(x)
}{
[v(x)]^2
}.
}
\]
\end{eigenschap}

Een geheugensteun is:

\[
\boxed{
\frac{
\text{afgeleide teller}\cdot\text{noemer}
-
\text{teller}\cdot\text{afgeleide noemer}
}{
\text{noemer}^2
}
}
\]

Let vooral op het \textbf{minteken}.


% ============================================================
% 5. EEN EERSTE RATIONALE FUNCTIE
% ============================================================

\FVDsection{Een eerste rationale functie afleiden}

\begin{voorbeeld}
Bepaal de afgeleide van
\[
f(x)=\frac{x^2+1}{x-2}.
\]

Eerst bepalen we het domein:
\[
\operatorname{dom}f=\mathbb{R}\setminus\{2\}.
\]

We kiezen
\[
u(x)=x^2+1,
\qquad
v(x)=x-2.
\]

Dan:
\[
u'(x)=2x,
\qquad
v'(x)=1.
\]

Met de quotiëntregel:
\[
f'(x)
=
\frac{2x(x-2)-(x^2+1)\cdot1}{(x-2)^2}.
\]

Vereenvoudigen geeft
\[
\begin{aligned}
f'(x)
&=
\frac{2x^2-4x-x^2-1}{(x-2)^2}\\
&=
\frac{x^2-4x-1}{(x-2)^2}.
\end{aligned}
\]

Dus
\[
\boxed{
f'(x)=\frac{x^2-4x-1}{(x-2)^2}
}
\qquad (x\neq2).
\]
\end{voorbeeld}

Merk op dat het domein niet plots verandert doordat we een afgeleide
berekenen. De oorspronkelijke functie is niet gedefinieerd voor $x=2$,
dus daar spreken we ook niet over haar afgeleide.


% ============================================================
% 6. EEN VASTE AANPAK
% ============================================================

\FVDsection{Een vaste aanpak voor de quotiëntregel}

Bij ingewikkeldere uitdrukkingen helpt een systematische aanpak.


Voor
\[
f(x)=\frac{u(x)}{v(x)}
\]
kun je als volgt werken:

\begin{enumerate}
  \item bepaal eerst waar de functie gedefinieerd is;
  \item herken teller $u$ en noemer $v$;
  \item bepaal afzonderlijk $u'$ en $v'$;
  \item schrijf de quotiëntregel op;
  \item vul zorgvuldig in;
  \item vereenvoudig pas daarna;
  \item controleer of je antwoord past bij het domein.
\end{enumerate}



\begin{voorbeeld}
Bepaal de afgeleide van
\[
f(x)=\frac{2x^3-x}{x^2+1}.
\]

Omdat
\[
x^2+1>0
\]
voor alle reële $x$, geldt
\[
\operatorname{dom}f=\mathbb{R}.
\]

Neem
\[
u(x)=2x^3-x,
\qquad
v(x)=x^2+1.
\]

Dan
\[
u'(x)=6x^2-1,
\qquad
v'(x)=2x.
\]

Dus
\[
f'(x)
=
\frac{
(6x^2-1)(x^2+1)-(2x^3-x)(2x)
}{
(x^2+1)^2
}.
\]

We werken de teller uit:
\[
\begin{aligned}
(6x^2-1)(x^2+1)
&=6x^4+5x^2-1,\\
(2x^3-x)(2x)
&=4x^4-2x^2.
\end{aligned}
\]

Daarom
\[
f'(x)
=
\frac{
2x^4+7x^2-1
}{
(x^2+1)^2
}.
\]

Dus
\[
\boxed{
f'(x)=
\frac{2x^4+7x^2-1}{(x^2+1)^2}.
}
\]
\end{voorbeeld}


% ============================================================
% 7. EERST HERSCHRIJVEN OF QUOTIËNTREGEL?
% ============================================================

\FVDsection{Eerst herschrijven of de quotiëntregel gebruiken?}

Niet elke breuk vraagt om de quotiëntregel.

Vergelijk de volgende functies:
\[
f(x)=\frac{3}{x^2}
\qquad\text{en}\qquad
g(x)=\frac{x^2+1}{x-2}.
\]

Voor $f$ kunnen we schrijven
\[
f(x)=3x^{-2}.
\]

De machtsregel geeft dan onmiddellijk
\[
f'(x)=-6x^{-3}
=
-\frac{6}{x^3}.
\]

De quotiëntregel gebruiken zou hier correct kunnen zijn, maar is
onnodig omslachtig.

Voor
\[
g(x)=\frac{x^2+1}{x-2}
\]
is zo'n eenvoudige herschrijving niet beschikbaar.
Daar is de quotiëntregel een natuurlijke keuze.


Kijk vóór het afleiden naar de structuur van de functie.

\[
\boxed{
\begin{array}{c}
\text{Kan ik eenvoudig herschrijven?}\\[1mm]
\downarrow\\[1mm]
\text{ja: herschrijf eerst}\\
\text{nee: kies een geschikte afleidingsregel}
\end{array}
}
\]



% ============================================================
% 8. VEREENVOUDIGEN IS GEEN BIJZAAK
% ============================================================

\FVDsection{Waarom vereenvoudigen belangrijk is}

Bij een afgeleide is een vereenvoudigde vorm vaak veel informatiever dan
de vorm die rechtstreeks uit de afleidingsregel ontstaat.

\begin{voorbeeld}
Neem
\[
f(x)=\frac{x^2}{x+1}.
\]

De quotiëntregel geeft
\[
f'(x)
=
\frac{2x(x+1)-x^2}{(x+1)^2}.
\]

Dat is correct, maar we kunnen de teller vereenvoudigen:
\[
2x(x+1)-x^2
=
2x^2+2x-x^2
=
x^2+2x.
\]

Factoriseren geeft:
\[
x^2+2x=x(x+2).
\]

Dus
\[
\boxed{
f'(x)=\frac{x(x+2)}{(x+1)^2}.
}
\]

Deze vorm is later veel nuttiger.

Omdat
\[
(x+1)^2>0
\qquad (x\neq-1),
\]
wordt het teken van $f'$ volledig bepaald door
\[
x(x+2).
\]

We kunnen daardoor veel gemakkelijker onderzoeken waar $f$ stijgt of
daalt.
\end{voorbeeld}

Dit is precies waarom algebra en analyse niet los van elkaar staan:
een goede algebraïsche vorm maakt de eigenschappen van een functie
zichtbaar.





% ============================================================

% 10. KIES JE STRATEGIE
% ============================================================

\FVDsection{Kies je strategie}

Bij het afleiden moet je steeds vaker eerst beslissen \emph{wat voor
functie je voor je hebt}.

\begin{center}
\[
\begin{array}{c|c}
\text{structuur} & \text{mogelijke aanpak}\\
\hline
u+v & \text{somregel}\\[1mm]
u-v & \text{verschilregel}\\[1mm]
c\,u & \text{constante-factorregel}\\[1mm]
uv & \text{productregel}\\[1mm]
\dfrac{u}{v} & \text{quotiëntregel}\\
\end{array}
\]
\end{center}

Maar deze tabel is geen automatisch stappenplan.

Bijvoorbeeld
\[
\frac{5}{x^3}
\]
ziet eruit als een quotiënt, maar
\[
5x^{-3}
\]
is veel gemakkelijker af te leiden.

De centrale vraag wordt dus:

\[
\boxed{
\text{Welke voorstelling van de functie maakt het afleiden het eenvoudigst?}
}
\]




\begin{oefening}[Productregel herkennen]{\oefB\ESS}

  Bepaal de afgeleide.
  
  \begin{question}
  
  \[
  f(x)=(x^2+1)(x-3).
  \]
  
  \begin{oplossing}
  We herkennen een product van
  \[
  u(x)=x^2+1
  \qquad\text{en}\qquad
  v(x)=x-3.
  \]
  
  Met de productregel:
  \[
  (uv)'=u'v+uv'.
  \]
  
  Omdat
  \[
  u'(x)=2x
  \qquad\text{en}\qquad
  v'(x)=1,
  \]
  krijgen we
  \[
  \begin{aligned}
  f'(x)
  &=2x(x-3)+(x^2+1)\cdot1\\
  &=2x^2-6x+x^2+1\\
  &=3x^2-6x+1.
  \end{aligned}
  \]
  
  Dus
  \[
  \boxed{f'(x)=3x^2-6x+1}.
  \]
  \end{oplossing}
  
  \end{question}
  
  
  \begin{question}
  
  \[
  g(x)=(2x^3-x)(x^2+4).
  \]
  
  \begin{oplossing}
  We nemen
  \[
  u(x)=2x^3-x,
  \qquad
  v(x)=x^2+4.
  \]
  
  Dan
  \[
  u'(x)=6x^2-1,
  \qquad
  v'(x)=2x.
  \]
  
  Met de productregel:
  \[
  \begin{aligned}
  g'(x)
  &=(6x^2-1)(x^2+4)+(2x^3-x)(2x)\\
  &=6x^4+23x^2-4+4x^4-2x^2\\
  &=10x^4+21x^2-4.
  \end{aligned}
  \]
  
  Dus
  \[
  \boxed{g'(x)=10x^4+21x^2-4}.
  \]
  \end{oplossing}
  
  \end{question}
  
  
  \begin{question}
  
  \[
  h(x)=x^2(3x-5).
  \]
  
  Bereken de afgeleide op twee manieren:
  
  \begin{enumerate}
    \item met de productregel;
    \item door eerst uit te werken.
  \end{enumerate}
  
  Vergelijk beide methoden.
  
  \begin{oplossing}
  \textbf{Methode 1: productregel}
  
  We nemen
  \[
  u(x)=x^2,
  \qquad
  v(x)=3x-5.
  \]
  
  Dan
  \[
  u'(x)=2x,
  \qquad
  v'(x)=3.
  \]
  
  Dus
  \[
  \begin{aligned}
  h'(x)
  &=2x(3x-5)+x^2\cdot3\\
  &=6x^2-10x+3x^2\\
  &=9x^2-10x.
  \end{aligned}
  \]
  
  \textbf{Methode 2: eerst uitwerken}
  
  We schrijven
  \[
  h(x)=3x^3-5x^2.
  \]
  
  Daaruit volgt rechtstreeks
  \[
  h'(x)=9x^2-10x.
  \]
  
  Beide methoden geven dus hetzelfde resultaat:
  \[
  \boxed{h'(x)=9x^2-10x}.
  \]
  
  Hier is eerst uitwerken iets korter. De productregel is correct,
  maar niet noodzakelijk de meest efficiënte methode.
  \end{oplossing}
  
  \end{question}
  
  \end{oefening}
  
  
  
  \begin{oefening}[Quotiëntregel]{\oefB\ESS}
  
  Bepaal telkens eerst het domein en vervolgens de afgeleide.
  
  \begin{question}
  
  \[
  f(x)=\frac{x^2+3}{x-1}.
  \]
  
  \begin{oplossing}
  De noemer mag niet nul zijn:
  \[
  x-1\neq0.
  \]
  
  Dus
  \[
  \boxed{\operatorname{dom}f=\mathbb{R}\setminus\{1\}}.
  \]
  
  We nemen
  \[
  u(x)=x^2+3,
  \qquad
  v(x)=x-1.
  \]
  
  Dan
  \[
  u'(x)=2x,
  \qquad
  v'(x)=1.
  \]
  
  Met de quotiëntregel:
  \[
  \begin{aligned}
  f'(x)
  &=
  \frac{2x(x-1)-(x^2+3)\cdot1}{(x-1)^2}\\
  &=
  \frac{2x^2-2x-x^2-3}{(x-1)^2}\\
  &=
  \frac{x^2-2x-3}{(x-1)^2}.
  \end{aligned}
  \]
  
  De teller kan worden ontbonden:
  \[
  x^2-2x-3=(x-3)(x+1).
  \]
  
  Dus
  \[
  \boxed{
  f'(x)=\frac{(x-3)(x+1)}{(x-1)^2}
  }
  \qquad (x\neq1).
  \]
  \end{oplossing}
  
  \end{question}
  
  
  \begin{question}
  
  \[
  g(x)=\frac{2x-1}{x^2+1}.
  \]
  
  \begin{oplossing}
  Omdat
  \[
  x^2+1>0
  \]
  voor alle reële $x$, geldt
  \[
  \boxed{\operatorname{dom}g=\mathbb{R}}.
  \]
  
  We nemen
  \[
  u(x)=2x-1,
  \qquad
  v(x)=x^2+1.
  \]
  
  Dan
  \[
  u'(x)=2,
  \qquad
  v'(x)=2x.
  \]
  
  Dus
  \[
  \begin{aligned}
  g'(x)
  &=
  \frac{2(x^2+1)-(2x-1)(2x)}
       {(x^2+1)^2}\\
  &=
  \frac{2x^2+2-4x^2+2x}
       {(x^2+1)^2}\\
  &=
  \frac{-2x^2+2x+2}
       {(x^2+1)^2}.
  \end{aligned}
  \]
  
  Dus
  \[
  \boxed{
  g'(x)=\frac{-2x^2+2x+2}{(x^2+1)^2}
  }.
  \]
  \end{oplossing}
  
  \end{question}
  
  
  \begin{question}
  
  \[
  h(x)=\frac{x^3+x}{x^2-4}.
  \]
  
  \begin{oplossing}
  De noemer mag niet nul zijn:
  \[
  x^2-4\neq0.
  \]
  
  Dus
  \[
  x\neq-2
  \qquad\text{en}\qquad
  x\neq2.
  \]
  
  Daarom
  \[
  \boxed{
  \operatorname{dom}h=\mathbb{R}\setminus\{-2,2\}
  }.
  \]
  
  We nemen
  \[
  u(x)=x^3+x,
  \qquad
  v(x)=x^2-4.
  \]
  
  Dan
  \[
  u'(x)=3x^2+1,
  \qquad
  v'(x)=2x.
  \]
  
  Met de quotiëntregel:
  \[
  h'(x)
  =
  \frac{
  (3x^2+1)(x^2-4)-(x^3+x)(2x)
  }{
  (x^2-4)^2
  }.
  \]
  
  We vereenvoudigen de teller:
  \[
  \begin{aligned}
  (3x^2+1)(x^2-4)
  &=3x^4-11x^2-4,\\
  (x^3+x)(2x)
  &=2x^4+2x^2.
  \end{aligned}
  \]
  
  Dus
  \[
  \begin{aligned}
  h'(x)
  &=
  \frac{
  3x^4-11x^2-4-(2x^4+2x^2)
  }{
  (x^2-4)^2
  }\\
  &=
  \frac{x^4-13x^2-4}{(x^2-4)^2}.
  \end{aligned}
  \]
  
  Daarom
  \[
  \boxed{
  h'(x)=\frac{x^4-13x^2-4}{(x^2-4)^2}
  }
  \qquad (x\neq\pm2).
  \]
  \end{oplossing}
  
  \end{question}
  
  \end{oefening}
  
  
  
  \begin{oefening}[Welke methode kies je?]{\oefU\ESS}
  
  Bepaal de afgeleide. Kies telkens eerst een geschikte strategie en
  motiveer kort waarom je die kiest.
  
  \begin{question}
  
  \[
  f(x)=\frac{4}{x^3}.
  \]
  
  \begin{oplossing}
  Hoewel $f$ als een quotiënt geschreven is, is de quotiëntregel hier
  onnodig omslachtig.
  
  We herschrijven:
  \[
  f(x)=4x^{-3},
  \qquad x\neq0.
  \]
  
  Met de machtsregel:
  \[
  f'(x)=4(-3)x^{-4}.
  \]
  
  Dus
  \[
  \boxed{
  f'(x)=-12x^{-4}=-\frac{12}{x^4}
  }
  \qquad (x\neq0).
  \]
  
  \textbf{Strategie:} eerst herschrijven en daarna de machtsregel gebruiken.
  \end{oplossing}
  
  \end{question}
  
  
  \begin{question}
  
  \[
  g(x)=x(x^2-5).
  \]
  
  \begin{oplossing}
  We kunnen de productregel gebruiken, maar eerst uitwerken is hier
  nog eenvoudiger:
  \[
  g(x)=x^3-5x.
  \]
  
  Dus
  \[
  \boxed{
  g'(x)=3x^2-5
  }.
  \]
  
  Met de productregel zouden we hetzelfde krijgen:
  \[
  g'(x)=1\cdot(x^2-5)+x\cdot2x
       =3x^2-5.
  \]
  
  \textbf{Strategie:} eerst uitwerken is hier het meest efficiënt.
  \end{oplossing}
  
  \end{question}
  
  
  \begin{question}
  
  \[
  h(x)=\frac{x^2+2}{x+1}.
  \]
  
  \begin{oplossing}
  De teller en noemer kunnen niet eenvoudig zo worden herschreven
  dat de breuk verdwijnt. De quotiëntregel is daarom een natuurlijke keuze.
  
  Het domein is
  \[
  \operatorname{dom}h=\mathbb{R}\setminus\{-1\}.
  \]
  
  Met
  \[
  u(x)=x^2+2,
  \qquad
  v(x)=x+1
  \]
  geldt
  \[
  u'(x)=2x,
  \qquad
  v'(x)=1.
  \]
  
  Dus
  \[
  \begin{aligned}
  h'(x)
  &=
  \frac{
  2x(x+1)-(x^2+2)
  }{
  (x+1)^2
  }\\
  &=
  \frac{x^2+2x-2}{(x+1)^2}.
  \end{aligned}
  \]
  
  Daarom
  \[
  \boxed{
  h'(x)=\frac{x^2+2x-2}{(x+1)^2}
  }
  \qquad (x\neq-1).
  \]
  
  \textbf{Strategie:} quotiëntregel.
  \end{oplossing}
  
  \end{question}
  
  
  \begin{question}
  
  \[
  p(x)=\frac{x^4-2x^2}{x^2}.
  \]
  
  \begin{oplossing}
  Eerst bepalen we het domein:
  \[
  x\neq0.
  \]
  
  Daarna vereenvoudigen we:
  \[
  p(x)
  =
  \frac{x^4}{x^2}
  -
  \frac{2x^2}{x^2}
  =
  x^2-2
  \qquad (x\neq0).
  \]
  
  Dus
  \[
  \boxed{
  p'(x)=2x
  }
  \qquad (x\neq0).
  \]
  
  \textbf{Strategie:} eerst vereenvoudigen. De quotiëntregel zou correct
  zijn, maar is hier onnodig lang.
  
  Merk op dat $x=0$ uitgesloten blijft, ook al komt die beperking
  niet meer voor in de vereenvoudigde formule.
  \end{oplossing}
  
  \end{question}
  
  \end{oefening}
  
  
  
  \begin{oefening}[Een fout analyseren]{\oefU\ESS}
  
  Een leerling berekent voor
  \[
  f(x)=\frac{x^2+1}{x-2}
  \]
  de afgeleide als volgt:
  \[
  f'(x)=\frac{2x}{1}=2x.
  \]
  
  \begin{question}
  
  Leg precies uit welke fout de leerling maakt.
  
  \begin{oplossing}
  De leerling heeft teller en noemer afzonderlijk afgeleid en daarna
  de resultaten gedeeld:
  \[
  \left(\frac{u}{v}\right)'
  \stackrel{\text{fout}}{=}
  \frac{u'}{v'}.
  \]
  
  Dat is geen geldige afleidingsregel.
  
  Voor een quotiënt geldt:
  \[
  \boxed{
  \left(\frac{u}{v}\right)'
  =
  \frac{u'v-uv'}{v^2}
  }.
  \]
  
  De verandering van teller én noemer moet dus samen in rekening
  worden gebracht.
  \end{oplossing}
  
  \end{question}
  
  
  \begin{question}
  
  Bereken de juiste afgeleide.
  
  \begin{oplossing}
  Het domein is
  \[
  \operatorname{dom}f=\mathbb{R}\setminus\{2\}.
  \]
  
  Neem
  \[
  u(x)=x^2+1,
  \qquad
  v(x)=x-2.
  \]
  
  Dan
  \[
  u'(x)=2x,
  \qquad
  v'(x)=1.
  \]
  
  Met de quotiëntregel:
  \[
  \begin{aligned}
  f'(x)
  &=
  \frac{
  2x(x-2)-(x^2+1)\cdot1
  }{
  (x-2)^2
  }\\
  &=
  \frac{
  2x^2-4x-x^2-1
  }{
  (x-2)^2
  }\\
  &=
  \frac{x^2-4x-1}{(x-2)^2}.
  \end{aligned}
  \]
  
  Dus
  \[
  \boxed{
  f'(x)=\frac{x^2-4x-1}{(x-2)^2}
  }
  \qquad (x\neq2).
  \]
  \end{oplossing}
  
  \end{question}
  
  
  \begin{question}
  
  Leg uit waarom de gebruikte foutieve regel niet de quotiëntregel kan zijn.
  
  \begin{oplossing}
  De quotiëntregel zegt niet dat teller en noemer afzonderlijk worden
  afgeleid.
  
  De correcte structuur is
  \[
  \frac{u'v-uv'}{v^2}.
  \]
  
  In het antwoord van de leerling ontbreken dus:
  
  \begin{itemize}
    \item de oorspronkelijke noemer $v$ bij de eerste term;
    \item de oorspronkelijke teller $u$ bij de tweede term;
    \item het verschil tussen beide producten;
    \item het kwadraat $v^2$ in de noemer.
  \end{itemize}
  
  De berekening
  \[
  \frac{u'}{v'}
  \]
  heeft daarom niets te maken met de quotiëntregel voor afgeleiden.
  
  Later zal een uitdrukking van de vorm
  \[
  \frac{u'}{v'}
  \]
  wel opduiken bij de regel van l'Hôpital, maar daar gaat het om
  \emph{limieten} en niet om de afgeleide van een quotiënt.
  \end{oplossing}
  
  \end{question}
  
  \end{oefening}




\end{document}